\begin{pastexam}[label={px:gw3-101-ds-06}]{關務三等資訊處理 101 年 第六題}
（申論）

每筆記錄是一對（關鍵值, 資料值）（key, data values），當 N 筆資料以下列資料結構儲存時：(1)2-3 樹（2-3 tree）
(2)AVL 樹（AVL tree） (3)最大堆（Max-Heap） (4)排序陣列（sorted array increasing order），試就搜尋（search key）、
刪除（delete key）、插入（insert key）、列印全部排序（print all nodes in order）及找最大值（find Max）等運算，比較其
時間複雜度。（20 分）\\
［註：請以表格列表呈現］
\end{pastexam}

\begin{pastsoln}
答案（最壞情況）：
\begin{center}
\begin{tabular}{lccccc}
\toprule
 & 搜尋 & 刪除 & 插入 & 依序列印 & 找最大值 \\
\midrule
(1) 2-3 樹 & $O(\log N)$ & $O(\log N)$ & $O(\log N)$ & $O(N)$ & $O(\log N)$ \\
(2) AVL 樹 & $O(\log N)$ & $O(\log N)$ & $O(\log N)$ & $O(N)$ & $O(\log N)$ \\
(3) 最大堆積 & $O(N)$ & $O(N)$ & $O(\log N)$ & $O(N \log N)$ & $O(1)$ \\
(4) 排序陣列 & $O(\log N)$ & $O(N)$ & $O(N)$ & $O(N)$ & $O(1)$ \\
\bottomrule
\end{tabular}
\end{center}

\textbf{(1) 2-3 樹}（資結書第~11.1~節）。 所有葉子在同一層，每個內部節點有 $2$ 或 $3$ 個子節點，高度介於
$\log_3 (N+1)$ 和 $\lg(N+1)$ 之間，是 $O(\log N)$。 搜尋從根往下，每層比較至多 $2$ 個鍵，$O(\log N)$。 插入找到葉子
放入，溢位就分裂、中間鍵往上送，每層最多分裂一次，$O(\log N)$。 刪除換成葉子上的鍵後刪除，不足就向兄弟借或與
兄弟合併，也是每層最多一次，$O(\log N)$。 依序列印就是中序走訪（每個節點的鍵和子樹交錯輸出），每個鍵輸出一次，
$O(N)$。 最大值在最右邊的葉子的最後一個鍵，從根一路往最右走，$O(\log N)$。

\textbf{(2) AVL 樹}（資結書第~10.2~節）。 高度 $< 1.44 \lg(N+2)$，是 $O(\log N)$。 搜尋 $O(\log N)$；插入最多一次
（單或雙）旋轉，刪除最多沿路徑每層旋轉一次，加上沿路更新高度，都是 $O(\log N)$。 中序走訪 $O(N)$ 依序列印。
最大值是一直往右子節點走到底，$O(\log N)$。 （2-3 樹、AVL 樹若另外維護一個指向最大節點的指標，找最大值可以變成
$O(1)$，但插入刪除時要順便更新。）

\textbf{(3) 最大堆積}（資結書第~8.1~節）。 堆積只保證父節點不小於子節點，兄弟之間、左右子樹之間沒有大小關係，
不能像搜尋樹那樣往一邊找，搜尋一個鍵最壞要看過全部，$O(N)$。 刪除任意一個鍵：先找到它 $O(N)$，再用最後一個元素
補位、上浮或下沉 $O(\log N)$，共 $O(N)$（若事先知道位置，或刪除的是最大值，只要 $O(\log N)$）。 插入放在最後
往上浮，$O(\log N)$。 依序列印要反覆取出最大值，$N$ 次各 $O(\log N)$，共 $O(N \log N)$，等於做一次堆積排序
（資結書第~14.6~節）。 不可能更快：建堆積只要 $O(N)$，如果堆積能在 $o(N \log N)$ 時間內依序列印，
合起來就是一個比 $\Omega(N \log N)$ 還快的比較排序，和排序的下界矛盾（資結書第~15.1~節）。
最大值就在根，$O(1)$。

\textbf{(4) 由小到大排序的陣列}。 搜尋用二分搜尋，$O(\log N)$（資結書第~17.1~節）。 插入先用二分搜尋找到位置
$O(\log N)$，但後面的元素要整段往後挪一格，最壞 $O(N)$；刪除同樣要把後面的元素往前挪，$O(N)$。 依序列印直接從頭
掃到尾，$O(N)$。 最大值是最後一個元素，$O(1)$。

\textbf{比較。} 需要頻繁搜尋、插入、刪除，又要能依序輸出時，平衡搜尋樹（2-3 樹、AVL 樹）搜尋、插入、刪除都是 $O(\log N)$，
依序列印 $O(N)$，最均衡；只關心「取最大值」並插入新資料時（優先佇列），最大堆積的 $O(1)$ 找最大、$O(\log N)$ 插入最划算，但不適合
搜尋；資料幾乎不變、主要是查詢時，排序陣列最省空間、搜尋也快，但插入刪除 $O(N)$。
\end{pastsoln}
